% ═══════════════════════════════════════════════════════════════
% MT-07d-ML: Musterlösung Minitest De compras (Wiederholung)
% Spanisch Klasse 7 - Sequenz 7d
% ═══════════════════════════════════════════════════════════════
% Dauer: 10-15 Minuten
% Punkte: 28 BE
% Inhalt: Lebensmittel, Zahlen, querer/poder, Einkaufsdialog
% ═══════════════════════════════════════════════════════════════

\documentclass[11pt, a4paper]{article}

\newif\ifswmodus
\swmodusfalse

\usepackage[utf8]{inputenc}
\usepackage[T1]{fontenc}
\usepackage[ngerman]{babel}
\usepackage{helvet}
\renewcommand{\familydefault}{\sfdefault}

\usepackage[margin=2cm]{geometry}
\usepackage{enumitem}
\usepackage{tabularx}
\usepackage{booktabs}

\usepackage{xcolor}
\usepackage{tcolorbox}
\tcbuselibrary{skins}

\usepackage{fancyhdr}
\usepackage{lastpage}
\usepackage{needspace}

% FARBEN
\definecolor{spanisch-color}{HTML}{c41e3a}
\definecolor{grau-color}{HTML}{888888}
\definecolor{hellgrau-color}{HTML}{f5f5f5}
\definecolor{orange-color}{HTML}{b35c00}
\definecolor{loesung-color}{HTML}{c62828}

\definecolor{sw-dark}{HTML}{000000}
\definecolor{sw-gray}{HTML}{666666}
\definecolor{sw-white}{HTML}{ffffff}

\ifswmodus
    \colorlet{spanisch}{sw-dark}
    \colorlet{grau}{sw-gray}
    \colorlet{hellgrau}{sw-white}
    \colorlet{orange}{sw-dark}
    \colorlet{loesung}{sw-dark}
\else
    \colorlet{spanisch}{spanisch-color}
    \colorlet{grau}{grau-color}
    \colorlet{hellgrau}{hellgrau-color}
    \colorlet{orange}{orange-color}
    \colorlet{loesung}{loesung-color}
\fi

\newcommand{\fachfarbe}{spanisch}

\newcommand{\thema}{Minitest: De compras -- MUSTERLÖSUNG}
\newcommand{\fach}{Spanisch}
\newcommand{\klasse}{7}
\newcommand{\maxpunkte}{28}
\newcommand{\zeit}{10-15 Minuten}
\newcommand{\datum}{\today}

\pagestyle{fancy}
\fancyhf{}
\renewcommand{\headrulewidth}{0.4pt}
\renewcommand{\footrulewidth}{0.4pt}

\lhead{\fach{} -- Klasse \klasse}
\chead{\textbf{MT-07d-ML}}
\rhead{\datum}

\lfoot{\zeit}
\cfoot{}
\rfoot{Seite \thepage\ von \pageref{LastPage}}

\newcommand{\punktebox}[1]{\hfill\fbox{\small \_/#1}}
\newcommand{\luecke}[1]{\rule{#1}{0.4pt}}
\newcommand{\lsg}[1]{\textcolor{loesung}{\textbf{#1}}}

\newtcolorbox{infobox}{
    colback=hellgrau,
    colframe=\fachfarbe,
    boxrule=1pt,
    arc=0pt,
    left=8pt, right=8pt, top=8pt, bottom=8pt
}

\newtcolorbox{hinweis}{
    colback=white,
    colframe=orange,
    boxrule=1pt,
    arc=0pt,
    left=5pt, right=5pt, top=5pt, bottom=5pt
}

\newenvironment{aufgabe}[1]{%
    \needspace{5\baselineskip}%
    \nopagebreak[4]%
    \vspace{10pt}
    \noindent\textbf{\textcolor{\fachfarbe}{Aufgabe #1}}
    \nopagebreak[4]%
    \vspace{5pt}
    \nopagebreak[4]%
    \begin{list}{}{\leftmargin=0pt}
    \item[]
}{%
    \end{list}
}

\newcommand{\es}[1]{\textcolor{\fachfarbe}{\textbf{#1}}}

\begin{document}

% ═══════════════════════════════════════════════════════════════
% KOPFBEREICH
% ═══════════════════════════════════════════════════════════════

\begin{center}
    {\Large\bfseries\textcolor{loesung}{\thema}}
\end{center}

\vspace{5pt}

\noindent
\begin{tabularx}{\textwidth}{|X|X|X|}
\hline
\textbf{Name:} \lsg{Musterlösung} &
\textbf{Klasse:} \klasse &
\textbf{Datum:} \luecke{2cm} \\
\hline
\end{tabularx}

\vspace{10pt}

\begin{infobox}
\textbf{Zeit:} \zeit{} | \textbf{Punkte:} \maxpunkte{} BE | \textbf{Hilfsmittel:} keine
\end{infobox}

\vspace{15pt}

% ═══════════════════════════════════════════════════════════════
% AUFGABE 1: Lebensmittel zuordnen (5 BE)
% ═══════════════════════════════════════════════════════════════

\begin{aufgabe}{1 -- Lebensmittel zuordnen}
Verbinde das spanische Wort mit der deutschen Bedeutung. \punktebox{5}
\end{aufgabe}

\begin{itemize}[nosep]
    \item \es{el queso} = \lsg{der Käse}
    \item \es{la manzana} = \lsg{der Apfel}
    \item \es{el pan} = \lsg{das Brot}
    \item \es{el pollo} = \lsg{das Hähnchen}
    \item \es{la leche} = \lsg{die Milch}
\end{itemize}

\vspace{15pt}

% ═══════════════════════════════════════════════════════════════
% AUFGABE 2: Zahlen schreiben (5 BE)
% ═══════════════════════════════════════════════════════════════

\begin{aufgabe}{2 -- Zahlen als Wort schreiben}
Schreibe die Zahl als spanisches Wort. \punktebox{5}
\end{aufgabe}

\noindent
a) 23 = \lsg{veintitrés}

\vspace{0.8em}

\noindent
b) 45 = \lsg{cuarenta y cinco}

\vspace{0.8em}

\noindent
c) 67 = \lsg{sesenta y siete}

\vspace{0.8em}

\noindent
d) 82 = \lsg{ochenta y dos}

\vspace{0.8em}

\noindent
e) 100 = \lsg{cien}

\vspace{15pt}

% ═══════════════════════════════════════════════════════════════
% AUFGABE 3: querer/poder Konjugation (6 BE)
% ═══════════════════════════════════════════════════════════════

\begin{aufgabe}{3 -- Konjugiere \es{querer} und \es{poder}}
Ergänze die richtige Form. \punktebox{6}
\end{aufgabe}

\noindent
a) Yo \lsg{quiero} comprar naranjas. \hfill (querer -- wollen)

\vspace{0.8em}

\noindent
b) ¿Tú \lsg{puedes} ayudarme? \hfill (poder -- können)

\vspace{0.8em}

\noindent
c) María \lsg{quiere} un kilo de tomates. \hfill (querer -- wollen)

\vspace{0.8em}

\noindent
d) Nosotros no \lsg{podemos} pagar con tarjeta. \hfill (poder -- können)

\vspace{0.8em}

\noindent
e) Ellos \lsg{quieren} ir al mercado. \hfill (querer -- wollen)

\vspace{0.8em}

\noindent
f) ¿Vosotros \lsg{podéis} hablar español? \hfill (poder -- können)

\vspace{15pt}

% ═══════════════════════════════════════════════════════════════
% AUFGABE 4: Einkaufsdialog vervollständigen (8 BE)
% ═══════════════════════════════════════════════════════════════

\begin{aufgabe}{4 -- Einkaufsdialog vervollständigen}
Ergänze den Dialog sinnvoll auf Spanisch. \punktebox{8}
\end{aufgabe}

\noindent
\textbf{Verkäufer:} ¡Buenos días! ¿Qué desea?

\vspace{0.8em}

\noindent
\textbf{Kunde:} \lsg{Quiero un kilo de naranjas.} \hfill (Ich möchte ein Kilo Orangen.)

\vspace{0.8em}

\noindent
\textbf{Verkäufer:} Aquí tiene. ¿Algo más?

\vspace{0.8em}

\noindent
\textbf{Kunde:} Sí, \lsg{¿puedo tener queso?} \hfill (Kann ich Käse haben?)

\vspace{0.8em}

\noindent
\textbf{Verkäufer:} Claro. ¿Cuánto quiere?

\vspace{0.8em}

\noindent
\textbf{Kunde:} 200 gramos. \lsg{¿Cuánto cuesta?} \hfill (Wie viel kostet das?)

\vspace{0.8em}

\noindent
\textbf{Verkäufer:} \lsg{Son seis euros con cincuenta.} \hfill (Das macht 6 Euro 50.)

\vspace{0.5em}

\textit{Bewertung: Je 2 BE pro korrekte Zeile (insgesamt 8 BE)}

\vspace{15pt}

% ═══════════════════════════════════════════════════════════════
% AUFGABE 5: Mengenangaben (4 BE)
% ═══════════════════════════════════════════════════════════════

\begin{aufgabe}{5 -- Mengenangaben}
Wähle die passende Mengenangabe. \punktebox{4}
\end{aufgabe}

\begin{hinweis}
\textbf{Mengenangaben:} un kilo de | un litro de | una botella de | un paquete de
\end{hinweis}

\vspace{0.5em}

\noindent
a) \lsg{un litro de} leche

\vspace{0.8em}

\noindent
b) \lsg{un kilo de} manzanas

\vspace{0.8em}

\noindent
c) \lsg{una botella de / un litro de} agua

\vspace{0.8em}

\noindent
d) \lsg{un paquete de} arroz

% ═══════════════════════════════════════════════════════════════
% PUNKTEÜBERSICHT
% ═══════════════════════════════════════════════════════════════

\vspace{25pt}
\hrule
\vspace{10pt}

\noindent
\begin{tabularx}{\textwidth}{|X|c|c|c|c|c||c|}
\hline
\textbf{Aufgabe} & 1 & 2 & 3 & 4 & 5 & \textbf{Gesamt} \\
\hline
\textbf{max. BE} & 5 & 5 & 6 & 8 & 4 & \textbf{28} \\
\hline
\textbf{erreicht} & & & & & & \\[0.8em]
\hline
\end{tabularx}

\vspace{15pt}

% ═══════════════════════════════════════════════════════════════
% 14-NP-SKALA
% ═══════════════════════════════════════════════════════════════

\noindent\textbf{14-NP-Bewertungsskala (Sek I):}

\vspace{0.5em}

\noindent
\begin{tabularx}{\textwidth}{|c|c|c|c|c|c|c|c|c|c|c|c|c|c|c|}
\hline
\textbf{NP} & 14 & 13 & 12 & 11 & 10 & 9 & 8 & 7 & 6 & 5 & 4 & 3 & 2 & 1 \\
\hline
\textbf{BE} & 28 & 27 & 25 & 24 & 22 & 21 & 19 & 17 & 15 & 13 & 11 & 9 & 8 & 6 \\
\hline
\end{tabularx}

\vspace{10pt}

\begin{flushright}
\textbf{Note:} \luecke{2cm}
\end{flushright}

\end{document}
